% GENERATED FILE. Edit canonical lessons, then run npm run generate:books.
% canonical-source-hash: fnv1a64:a4f38806639590a2
% canonical-lessons: BN-C80-deri, BN-C80-age, BN-C80-pore, BN-C80-somoy, BN-C80-dupur

\chapter{Lateness, Before, After, Time, Midday}
\label{ch:bn-lateness-before-after-time-midday}
\hlchaptermodality{80}

\begin{chapteropening}
\textbf{By the end of this chapter:} \emph{I can say lateness, before, after, later, time and midday.}

Complete the last lesson of chapter 80: I can say lateness, before, after, later, time and midday.
\end{chapteropening}

\section[\texorpdfstring{\bn{দেরি} (deri)}{দেরি (deri)}]{\bn{দেরি} (deri) — lateness; late}

\label{lesson:BN-C80-deri}



\begin{quote}
Before the new one: say the Bengali for now, then the Bengali for always.
\end{quote}

\subsection*{You'll want to know: \bn{দেরি}}
\textbf{\bn{দেরি}} — \emph{deri} — \textquotedblleft{}lateness; late\textquotedblright{}.

\begin{cousinweb}
\textbf{\bn{দেরি}} (\emph{deri}) is a loan, usually traced to Persian \emph{dēr}, \textquotedblleft{}late; a long time\textquotedblright{}. Bengali adds the ending \emph{-i}.
\end{cousinweb}

\begin{grammarlens}[title={What you've built}]
One more word for this chapter.
\end{grammarlens}

\subsection*{Your turn}
\begin{itemize}
\raggedright
  \item \emph{Say it:} \emph{deri}
  \item \emph{Say it:} \emph{deri}, once more
  \item \emph{Say it:} say \emph{deri}
  \item \emph{Recall:} say the Bengali for now, then the Bengali for always, then say \emph{deri} again
\end{itemize}

\begin{tcolorbox}[breakable,colback=teal!4,colframe=teal!35!black,title={Before you move on}]
What does \bn{দেরি} mean? (\textquotedblleft{}Lateness; late\textquotedblright{}.) Say it once more.
\end{tcolorbox}
\section[\texorpdfstring{\bn{আগে} (āge)}{আগে (āge)}]{\bn{আগে} (āge) — before}

\label{lesson:BN-C80-age}



\begin{quote}
Before the new one: say the Bengali for always, then the Bengali for lateness; late.
\end{quote}

\subsection*{You'll want to know: \bn{আগে}}
\textbf{\bn{আগে}} — \emph{āge} — \textquotedblleft{}before\textquotedblright{}.

\begin{cousinweb}
\textbf{\bn{আগে}} (\emph{āge}) is a tadbhava, a word inherited from Sanskrit and reshaped by sound change. It comes from \emph{agra}, \textquotedblleft{}top, front\textquotedblright{} (CDIAL 68). CDIAL lists \emph{āge} under the locative \emph{agre}, \textquotedblleft{}in front\textquotedblright{}.
\end{cousinweb}

\begin{grammarlens}[title={What you've built}]
One more word for this chapter.
\end{grammarlens}

\subsection*{Your turn}
\begin{itemize}
\raggedright
  \item \emph{Say it:} \emph{āge}
  \item \emph{Say it:} \emph{āge}, once more
  \item \emph{Say it:} say \emph{āge}
  \item \emph{Recall:} say the Bengali for always, then the Bengali for lateness; late, then say \emph{āge} again
\end{itemize}

\begin{tcolorbox}[breakable,colback=teal!4,colframe=teal!35!black,title={Before you move on}]
What does \bn{আগে} mean? (\textquotedblleft{}Before\textquotedblright{}.) Say it once more.
\end{tcolorbox}
\section[\texorpdfstring{\bn{পরে} (pôre)}{পরে (pôre)}]{\bn{পরে} (pôre) — after, later}

\label{lesson:BN-C80-pore}



\begin{quote}
Before the new one: say the Bengali for lateness; late, then the Bengali for before.
\end{quote}

\subsection*{You'll want to know: \bn{পরে}}
\textbf{\bn{পরে}} — \emph{pôre} — \textquotedblleft{}after, later\textquotedblright{}.

\begin{cousinweb}
\textbf{\bn{পরে}} (\emph{pôre}) is \emph{pôr}, \textquotedblleft{}later\textquotedblright{}, with the ending \emph{-e}. It goes back to Sanskrit \emph{para}, \textquotedblleft{}further, later\textquotedblright{}, whose \emph{param} means \textquotedblleft{}after\textquotedblright{}; CDIAL has that entry but lists no Bengali form under it (CDIAL 7793).
\end{cousinweb}

\begin{grammarlens}[title={What you've built}]
One more word for this chapter.
\end{grammarlens}

\subsection*{Your turn}
\begin{itemize}
\raggedright
  \item \emph{Say it:} \emph{pôre}
  \item \emph{Say it:} \emph{pôre}, once more
  \item \emph{Say it:} say \emph{pôre}
  \item \emph{Recall:} say the Bengali for lateness; late, then the Bengali for before, then say \emph{pôre} again
\end{itemize}

\begin{tcolorbox}[breakable,colback=teal!4,colframe=teal!35!black,title={Before you move on}]
What does \bn{পরে} mean? (\textquotedblleft{}After, later\textquotedblright{}.) Say it once more.
\end{tcolorbox}
\section[\texorpdfstring{\bn{সময়} (sômôy)}{সময় (sômôy)}]{\bn{সময়} (sômôy) — time}

\label{lesson:BN-C80-somoy}



\begin{quote}
Before the new one: say the Bengali for before, then the Bengali for after.
\end{quote}

\subsection*{You'll want to know: \bn{সময়}}
\textbf{\bn{সময়}} — \emph{sômôy} — \textquotedblleft{}time\textquotedblright{}.

\begin{cousinweb}
\textbf{\bn{সময়}} (\emph{sômôy}) is Sanskrit \emph{samaya}, \textquotedblleft{}meeting, agreement; time\textquotedblright{}, taken over whole as a learned loan, a tatsama.
\end{cousinweb}

\begin{grammarlens}[title={What you've built}]
One more word for this chapter.
\end{grammarlens}

\subsection*{Your turn}
\begin{itemize}
\raggedright
  \item \emph{Say it:} \emph{sômôy}
  \item \emph{Say it:} \emph{sômôy}, once more
  \item \emph{Say it:} say \emph{sômôy}
  \item \emph{Recall:} say the Bengali for before, then the Bengali for after, then say \emph{sômôy} again
\end{itemize}

\begin{tcolorbox}[breakable,colback=teal!4,colframe=teal!35!black,title={Before you move on}]
What does \bn{সময়} mean? (\textquotedblleft{}Time\textquotedblright{}.) Say it once more.
\end{tcolorbox}
\section[\texorpdfstring{\bn{দুপুর} (dupur)}{দুপুর (dupur)}]{\bn{দুপুর} (dupur) — midday}

\label{lesson:BN-C80-dupur}



\begin{quote}
Before the new one: say the Bengali for after, then the Bengali for time.
\end{quote}

\subsection*{You'll want to know: \bn{দুপুর}}
\textbf{\bn{দুপুর}} — \emph{dupur} — \textquotedblleft{}midday\textquotedblright{}.

\begin{cousinweb}
\textbf{\bn{দুপুর}} (\emph{dupur}) is usually explained as Sanskrit \emph{dviprahara}, \textquotedblleft{}the second watch\textquotedblright{} of the day: \emph{dvi}, \textquotedblleft{}two\textquotedblright{}, + \emph{prahara}, \textquotedblleft{}a watch of about three hours\textquotedblright{}, which CDIAL lists as Bengali \emph{pahar} (CDIAL 8900). CDIAL has no entry for the compound, and asks whether \emph{prahara} is a Sanskrit form of Persian \emph{pahra}, \textquotedblleft{}watch\textquotedblright{}.
\end{cousinweb}

\begin{grammarlens}[title={What you've built}]
One more word for this chapter.
\end{grammarlens}

\subsection*{Your turn}
\begin{itemize}
\raggedright
  \item \emph{Say it:} \emph{dupur}
  \item \emph{Say it:} \emph{dupur}, once more
  \item \emph{Say it:} say \emph{dupur}
  \item \emph{Recall:} say the Bengali for after, then the Bengali for time, then say \emph{dupur} again
\end{itemize}

\begin{tcolorbox}[breakable,colback=teal!4,colframe=teal!35!black,title={Before you move on}]
What does \bn{দুপুর} mean? (\textquotedblleft{}Midday\textquotedblright{}.) Say it once more.
\end{tcolorbox}
